% Additional foundational concepts: induction, convexity, and spectral algebra.

\concept{F35}{Mathematical induction proves a statement at every step}{F01, F02, F08}
\paragraph{Foundation and definition.}
Induction is a proof principle for statements indexed by the
nonnegative integers. It requires a base case and an implication from
each established case to the next. It can be taken as an axiom of the
natural numbers; equivalently, it follows from their well-ordering
principle, which says every nonempty set of nonnegative integers has a
least member. It is not an empirical extrapolation from several examples.
\paragraph{Symbols and meaning.}
Let $P(q)$ be a proposition indexed by the integer $q\ge0$. The base
case is $P(0)$. The induction step is
$\forall q\ge0:\ P(q)\Rightarrow P(q+1)$. Together they imply
$P(q)$ for every nonnegative integer $q$. For a finite horizon
$Q\ge1$, the step need only hold for integers $0\le q<Q$ to conclude
all cases $0\le q\le Q$.
\paragraph{Derivation or example.}
Suppose a counterexample exists despite the base and step. Its least
index $q_*$ cannot be zero because the base holds. Thus $q_*-1$ is a
nonnegative earlier index, where the claim holds by minimality. The step
then proves $P(q_*)$, a contradiction. For a state example, define
$\mathcal R_0=\{x_0\}$ and
$\mathcal R_{q+1}=\{T(x,u):x\in\mathcal R_q,u\in\mathcal U_{q+1}\}$.
Here states lie in a set $\mathcal X$, inputs lie in a set
$\mathcal U$, $x_0\in\mathcal X$ is the initial state,
$T:\mathcal X\times\mathcal U\to\mathcal X$ is a transition function,
and $\mathcal U_{q+1}\subseteq\mathcal U$ permits next inputs. Membership of each
actual input in its permitted set proves actual-state membership by
induction.
\paragraph{Why ECDR needs it.}
A finite-horizon reuse certificate must retain the effects of every
previous accepted wave. Induction connects trusted initial state and
valid transitions to membership in the certified reachable sets.
\paragraph{Conditions and connections.}
The step must cover every permitted previous state and next input.
Checking only one observed path proves no universal step. After a failed
input gate, continuing requires a new valid base or another proved rule.

\concept{F36}{Convex sets and convex combinations}{F03, F11, F16}
\paragraph{Foundation and definition.}
A set is convex when the full segment joining any two of its members
also belongs to it. This is a geometric definition, without probabilities
or class labels. A convex combination uses nonnegative coefficients
summing to one; unrestricted real coefficients instead define linear
combinations, and unrestricted coefficients summing to one define affine
combinations.
\paragraph{Symbols and meaning.}
Let $\mathcal S\subseteq\R^d$, where integer $d\ge1$ counts
coordinates. For $x,y\in\mathcal S$ and scalar $t\in[0,1]$,
convexity requires $tx+(1-t)y\in\mathcal S$. More generally, a convex
combination of $r\ge1$ points $x_1,\ldots,x_r$ uses coefficients
$\alpha_i\ge0$ satisfying $\sum_{i=1}^{r}\alpha_i=1$, and is
$\sum_i\alpha_i x_i$. The index $i$ selects a point.
\paragraph{Derivation or example.}
For a box $[\ell,u]\subseteq\R^d$ and members $x,y$, multiply each
coordinate inequality by $t$ or $1-t$ and add:
\[
 \ell_j\le tx_j+(1-t)y_j\le u_j,
 \qquad j=1,\ldots,d.
\]
Thus boxes are convex. Norm balls are convex because the triangle
inequality gives
$\norm{tx+(1-t)y}\le t\norm{x}+(1-t)\norm{y}$ for a ball centered
at zero. In contrast, $\{-1,1\}\subset\R$ excludes the midpoint
zero and is not convex. The convex hull of supplied points is the set
of all their finite convex combinations.
\paragraph{Why ECDR needs it.}
A scalar function $f:\mathcal S\to\R$ is convex when
$f(tx+(1-t)y)\le tf(x)+(1-t)f(y)$. ReLU satisfies this inequality,
so its endpoint chord upper-bounds its graph. Convex relaxations use
such enclosures. A convex input domain alone says nothing about a CNN's
decision throughout that domain.
\paragraph{Conditions and connections.}
Affine images and intersections preserve convexity; arbitrary unions
need not. A union of certified regions can be safe even when nonconvex,
provided each accepted query has a valid regional witness. Replacing a
union by its convex hull requires a new certificate on added points.

\concept{F37}{Eigenvectors and eigenvalues identify preserved directions}{F12, F13, F15}
\paragraph{Foundation and definition.}
An eigenvector of a square matrix is a nonzero vector whose direction
is preserved by the matrix, although its magnitude or sign may change.
The associated scaling factor is its eigenvalue. The zero vector is
excluded because it satisfies the equation for every scalar and would
identify no direction.
\paragraph{Symbols and meaning.}
Let $A\in\R^{d\times d}$ be a square real matrix, with integer
$d\ge1$. A real eigenpair consists of scalar $\lambda\in\R$ and
nonzero column $v\in\R^d$ satisfying
\[
 Av=\lambda v.
\]
The nonzero scalar multiples of $v$ describe the same direction. A real
matrix need not have real eigenvalues or a complete eigenvector basis.
Complex eigenpairs require extending the scalar field and are outside
this page's real-coordinate treatment.
\paragraph{Derivation or example.}
For $A=\begin{pmatrix}2&1\\1&2\end{pmatrix}$,
$v_1=(1,1)^\top$ satisfies $Av_1=3v_1$, while
$v_2=(1,-1)^\top$ satisfies $Av_2=v_2$. Thus the sum direction is
amplified by three and the difference direction by one.
Equivalently, $(A-\lambda I_d)v=0$, where $I_d$ is the identity
matrix. Thus $A-\lambda I_d$ sends a nonzero vector to zero and
cannot have an inverse: an inverse would send zero back to both zero
and $v$. The eigenvalue identifies precisely this degeneracy.
\paragraph{Why ECDR needs it.}
PCA uses eigenvectors of a covariance matrix to describe directions of
variation. The eigenvalues quantify variance along unit eigenvectors.
Large variance along a direction does not establish that a classifier's
margin is insensitive in that direction.
\paragraph{Conditions and connections.}
Generic square matrices can be defective or have complex eigenvalues.
Real symmetric matrices have the stronger orthonormal-basis guarantee
in F38. Repeated eigenvalues can make individual eigenvectors
nonunique, even while their eigenspace is well defined.

\concept{F38}{Symmetry, positive semidefiniteness, and the spectral theorem}{F18, F19, F27, F37}
\paragraph{Foundation and definition.}
A real square matrix is symmetric when transposition leaves it
unchanged. A symmetric matrix is positive semidefinite, abbreviated PSD,
when every quadratic form it produces is nonnegative. PSD does not mean
that every matrix entry is nonnegative.
\paragraph{Symbols and meaning.}
Let $A\in\R^{d\times d}$ satisfy $A^\top=A$, with integer
$d\ge1$. The PSD condition is $x^\top Ax\ge0$ for every
$x\in\R^d$. The finite-dimensional real spectral theorem states
\[
 A=V\Lambda V^\top,\qquad V^\top V=I_d,
 \qquad\Lambda=\diag(\lambda_1,\ldots,\lambda_d).
\]
Columns of $V\in\R^{d\times d}$ form an orthonormal eigenbasis;
$\lambda_i\in\R$ are eigenvalues, and $I_d$ is the identity.
This is an established theorem, rather than an additional axiom.
\paragraph{Derivation or example.}
A proof sketch maximizes $x^\top Ax$ on the compact unit sphere.
Its first-order condition produces a unit eigenvector. Symmetry makes
that vector's orthogonal complement invariant under $A$; induction on
dimension constructs the remaining eigenvectors. This sketches the
argument rather than supplying every optimization detail.
For observations $z_1,\ldots,z_N\in\R^d$ centered to have zero
mean, with integer $N\ge2$, the covariance matrix
$C=(N-1)^{-1}\sum_{i=1}^{N}z_i z_i^\top$ is symmetric and
\[
 x^\top Cx=\frac1{N-1}\sum_{i=1}^{N}(z_i^\top x)^2\ge0.
\]
For a unit eigenvector $v$, its eigenvalue is $v^\top Cv\ge0$.
\paragraph{Why ECDR needs it.}
The theorem justifies orthonormal PCA directions and nonnegative
explained-variance contributions. It also distinguishes a covariance
eigenvalue from a signed classifier weight.
\paragraph{Conditions and connections.}
For symmetric $A$, PSD is equivalent to all eigenvalues being
nonnegative. Zero eigenvalues are allowed and identify directions with
zero quadratic form. Finite-precision eigensolvers need numerical
checks; a tiny negative computed covariance eigenvalue may reflect
rounding rather than a negative mathematical variance.

\concept{F39}{Singular values connect PCA and operator amplification}{F15, F21, F38}
\paragraph{Foundation and definition.}
The singular value decomposition, or SVD, expresses any finite real
matrix through orthogonal input and output coordinates and nonnegative
axis scalings. Unlike eigenvalue decomposition, it applies to rectangular
matrices. Its existence is an established finite-dimensional theorem;
it follows by applying the spectral theorem to $A^\top A$ and
completing orthonormal bases for null spaces.
\paragraph{Symbols and meaning.}
For $A\in\R^{m\times n}$, with integers $m,n\ge1$, write
\[
 A=U\Sigma V^\top,\qquad U^\top U=I_m,\quad V^\top V=I_n.
\]
Here $U\in\R^{m\times m}$ and $V\in\R^{n\times n}$ are
orthogonal, $I_m,I_n$ are identity matrices, and
$\Sigma\in\R^{m\times n}$ is rectangular diagonal. Its singular
values satisfy $\sigma_1\ge\cdots\ge\sigma_p\ge0$, where
$p=\min(m,n)$. Column $v_i$ of $V$ is a right singular direction.
\paragraph{Derivation or example.}
For any input $x\in\R^n$, put $y=V^\top x$. Orthogonality gives
\[
 \norm{Ax}_2^2=\sum_{i=1}^{p}\sigma_i^2y_i^2
 \le\sigma_1^2\norm{x}_2^2.
\]
Taking $x=v_1$ attains the bound, so the induced Euclidean operator
norm is exactly $\norm{A}_{2\to2}=\sigma_1$.
For centered observations arranged as rows of
$X\in\R^{N\times d}$, with integer $N\ge2$, covariance is
$C=X^\top X/(N-1)$. Substituting $X=U_X\Sigma_XV_X^\top$ shows
that columns of $V_X$ are PCA directions; the nonzero covariance
eigenvalues are the squared singular values divided by $N-1$.
\paragraph{Why ECDR needs it.}
One decomposition connects empirical variation geometry and worst-case
linear amplification, while keeping their meanings distinct: variance
describes fitted observations, whereas an operator norm covers every
input direction.
\paragraph{Conditions and connections.}
Repeated singular values permit nonunique directions. A truncated SVD
describes an approximation, not a universal residual bound. Numerical
estimates of the largest singular value need a certified error allowance
before they serve as upper amplification bounds in a proof.

